\chapter{Drift correction sketch}
\hypertarget{md__sketch_examples_2_weight_reading_with_recalibration_mode_switching_2_r_e_a_d_m_e}{}\label{md__sketch_examples_2_weight_reading_with_recalibration_mode_switching_2_r_e_a_d_m_e}\index{Drift correction sketch@{Drift correction sketch}}
\label{md__sketch_examples_2_weight_reading_with_recalibration_mode_switching_2_r_e_a_d_m_e_autotoc_md0}%
\Hypertarget{md__sketch_examples_2_weight_reading_with_recalibration_mode_switching_2_r_e_a_d_m_e_autotoc_md0}%
 This Arduino sketch implements the drift correction strategy that consists of taking weight measurements when the scale is empty as the new reference/offset. This strategy is detailled in \href{https://www.overleaf.com/read/vvxvjbdjmgmg}{\texttt{ Lab notes}} at the end of the document.

In this sketch, the start-\/up is not done with easy\+\_\+start\+\_\+with\+\_\+params(). Functions that allow the taring of all the Load\+Cells are used.

Also, the controller.\+get\+\_\+weight\+\_\+with\+\_\+auto\+\_\+recalibration() function is used to perform a reading.

The code allows to switch from automatic start-\/up to manual calibration start-\/up (see sketch for more details) 